\documentclass[10pt]{article}
\usepackage[letterpaper,margin=0.75in,top=0.6in,bottom=0.6in]{geometry}
\usepackage[T1]{fontenc}
\usepackage{mathptmx}
\usepackage{booktabs}
\usepackage{enumitem}
\usepackage{titlesec}
\usepackage{microtype}

\pagestyle{empty}
\setlength{\parindent}{0pt}
\setlength{\parskip}{3pt}
\setlist[itemize]{leftmargin=1.2em,itemsep=1pt,topsep=1pt,parsep=0pt}
\titleformat{\section}{\normalsize\bfseries}{}{0pt}{}
\titlespacing*{\section}{0pt}{6pt}{2pt}

\begin{document}

\begin{center}
  {\large\bfseries Status Update: Belief-Space Planning for the Card Game Fish (Literature)}\\[2pt]
  [Group member names] \quad|\quad November 13, 2026
\end{center}

\section*{Problem}
Fish is a six-player, two-team card game of hidden information: 54 cards form nine sets of six,
players ask opponents for specific cards, and a team scores a set only by declaring exactly which
teammate holds each of its cards (a wrong declaration gives the set to the other team). Every ask,
its outcome and all hand sizes are public, so good play is mostly inference. Our goal is an agent
that (i) maintains a probabilistic belief over where every hidden card is and (ii) plans with that
belief, and that beats rule-based baselines. Reward is sets won minus sets lost per game
(range $-9$ to $+9$).

\section*{What we have accomplished}
\begin{itemize}
  \item \textbf{Engine.} Full rules with seeded deals; 10,000 random games complete with
    card-conservation checks. Rule ambiguities are resolved and documented.
  \item \textbf{Hard-constraint knowledge.} A public constraint store (who can hold which card,
    hand sizes, minimum cards per set) with propagation to a fixed point. In 3,300 simulated games
    the true deal never contradicted any player's view.
  \item \textbf{Particle belief.} A sampler over deals consistent with the constraints
    (constructive pass, importance resampling, MCMC swaps). On 12 endgames small enough to
    enumerate, sampled marginals are within 0.03 of exact. Speed is about 20,000 samples per second.
  \item \textbf{Agents.} Random, RuleBased, GreedyBelief (ask the likeliest card, declare above a
    confidence threshold), determinized Rollout, and single-observer ISMCTS. All see only their own
    hand and public information. Mean decision time is about 0.2\,s for the planners.
  \item \textbf{Evaluation harness.} Parallel arena with paired deals (teams swapped), CSV logging
    and plotting scripts; 70 automated tests pass.
\end{itemize}

\textbf{Results so far} (mean set differential per game, row agent minus column agent, 400 games
per pairing, 95\% confidence intervals):

\begin{center}
\small
\begin{tabular}{@{}lcccc@{}}
\toprule
 & \textbf{vs Random} & \textbf{vs RuleBased} & \textbf{vs GreedyBelief} & \textbf{vs Rollout} \\
\midrule
\textbf{RuleBased}    & $+7.36 \pm 0.18$ &                  &                  &                  \\
\textbf{GreedyBelief} & $+7.73 \pm 0.15$ & $+2.27 \pm 0.27$ &                  &                  \\
\textbf{Rollout}      & $+7.15 \pm 0.18$ & $+2.37 \pm 0.28$ & $+0.36 \pm 0.23$ &                  \\
\textbf{ISMCTS}       & $+7.18 \pm 0.18$ & $+2.50 \pm 0.25$ & $+0.47 \pm 0.27$ & $+0.30 \pm 0.26$ \\
\bottomrule
\end{tabular}
\end{center}

\begin{itemize}
  \item The main goal is met: every belief-based agent beats RuleBased by 2.3 to 2.5 sets per game.
    Planning adds a smaller but statistically significant edge over the greedy belief agent
    (about 0.4 sets per game).
  \item GreedyBelief's declare threshold has a plateau: strength is flat from 0.8 to 0.99 and falls
    off below about 0.6 (a threshold of 0.3 loses 1.5 sets per game to a threshold of 0.9).
  \item The belief is well calibrated overall (mean Brier 0.549 against 0.594 for a
    uniform-over-possible-holders baseline) but under-confident in the 0.4 to 0.8 range, because it
    ignores what a player's choice of ask says about their hand. Reweighting particles by a simple
    ask-likelihood model removes most of that gap (Brier 0.541); it is implemented but not yet
    switched on in the agents.
\end{itemize}

\section*{Open issues and where we would like feedback}
\begin{itemize}
  \item \textbf{Planning gain is modest.} Rollouts use a cheap approximation of the greedy policy
    for the simulated players and stop after 12 plies with a heuristic leaf value; deeper rollouts
    did not help in tuning. Is a learned or better-fitted leaf evaluation a reasonable next step, or
    should we invest in the opponent model instead?
  \item \textbf{Baselines.} Our strongest non-planning baseline is our own greedy agent.
    Suggestions for a stronger hand-written baseline would be welcome.
  \item \textbf{Not finished:} the search-budget versus win-rate curve was interrupted and must be
    re-run; the round-robin used 400 games per pairing instead of the planned 1,000.
\end{itemize}

\section*{Revised timeline}
\begin{tabular}{@{}lp{5.7in}@{}}
\textbf{Week of Nov 16} & Re-run the budget curve; round-robin at 1,000 games per pairing; evaluate
  ask-likelihood reweighting inside the agents. \\
\textbf{Week of Nov 23} & Improve the planners: better leaf evaluation, common random numbers
  across candidates, tuned opponent model. \\
\textbf{Week of Nov 30} & Final experiments with fixed seeds; ablations (belief quality versus
  playing strength); figures. \\
\textbf{Week of Dec 7}  & Final report and code clean-up. \\
\end{tabular}

\end{document}
